\section{Methodology}
\label{sec:method}

Our main question is whether instruct LLMs prefer early answer positions, by how much this changes measured accuracy, and how the effect depends on option count and question difficulty.
We fixed the hypotheses and the analysis plan before running any model.

\para{Hypotheses.}
We test five claims:
(H1) models over-select early positions (A/B);
(H2) on 4-option items this inflates accuracy by 3--8 points;
(H3) the effect is largest on medium-difficulty questions;
(H4) true/false shows $<$2 points and 5-option shows $>$10 points;
(H5) near the decision boundary, up to 15\% of apparent accuracy is positional.

\subsection{Within-question rotation design}
\label{sec:design}

For a question with $k$ options, we show all $k$ \emph{cyclic rotations} of its option list.
The correct answer (the \emph{key}) therefore occupies every position exactly once per question, and an unbiased model chooses each position $1/k$ of the time.
Rotation 0 of the native-$k$ item is the order in which the benchmark was published.
We choose rotations over full shuffles because they balance key position with only $k$ prompts per question; the cost is that they cover $k$ of the $k!$ possible orders.

To vary the number of options without changing content, we build reduced-$k$ versions of each question that keep the key plus a seeded, nested sample of the same question's distractors.
\Tabref{tab:conditions} lists all conditions.
In total each model answers 42{,}594 prompts, for 298{,}158 scored prompts overall.

\begin{table}[t]
    \centering
    \caption{Conditions. Each question appears once per cyclic rotation of its options, so the key is balanced across positions in every condition.}
    \label{tab:conditions}
    \small
    \begin{tabular}{@{}llr@{}}
        \toprule
        \textbf{Condition} & \textbf{Options shown} & \textbf{Prompts per model} \\
        \midrule
        \mmlu $k{=}4$, $k{=}3$, $k{=}2$ & A--D, A--C, A--B & 6{,}000 / 4{,}500 / 3{,}000 \\
        \mmlu $k{=}4$, numeric labels & 1--4 (label ablation) & 6{,}000 \\
        \csqashort $k{=}5$, $4$, $3$, $2$ & A--E down to A--B & 6{,}105 / 4{,}884 / 3{,}663 / 2{,}442 \\
        \boolq, A/B letters & ``A. True / B. False'' and the reverse & 3{,}000 \\
        \boolq, True/False words & ``True or False'' vs.\ ``False or True'' & 3{,}000 \\
        \bottomrule
    \end{tabular}
\end{table}

\subsection{Data, models, and prompts}

\para{Data.}
We use a sample of 1{,}500 \mmlu test questions stratified over the 57 subjects (seed 42), all 1{,}221 \csqa validation questions, and 1{,}500 \boolq validation questions with passages under 1{,}500 characters (seed 42).

\para{Models.}
We run seven open-weight instruct models locally: Qwen2.5 at 1.5B, 7B, and 14B, Qwen3-4B-Instruct-2507, Qwen3-8B with thinking disabled, Llama-3.1-8B-Instruct, and Phi-3-mini-4k-instruct.
We refer to them by the tags used in the tables (\eg qwen2.5-7b).
All models run in bfloat16 with HuggingFace \texttt{transformers}; \secref{app:setup} gives repository names and hardware.
We planned to include hosted frontier models but could not, because no API access was available; our conclusions therefore cover open-weight models of 14B parameters or fewer.

\para{Prompts.}
Prompts are zero-shot, with no system message, in each model's own chat template.
A lettered item lists the question, the labeled options, and the instruction ``Answer with the letter of the correct option only.''
\boolq items show a passage and the statement ``The answer to the question `\ldots?' is yes.''
In the word condition the model answers ``True'' or ``False'', and we swap the order in which the instruction mentions the two words; in the letter condition the two answers are lettered options A and B, shown in both orders.
\Secref{app:prompts} gives the templates verbatim.

\para{Read-out.}
We take the log-probability of each label at the first assistant token, summed over every vocabulary token that decodes to that label with or without leading whitespace.
The prediction is the highest-scoring label.
There is no sampling, so uncertainty comes from the choice of questions and not from generation.

\subsection{Metrics}
\label{sec:metrics}

Let $\acc_p$ be accuracy when the key is at position $p$, and let $\bal = \frac{1}{k}\sum_p \acc_p$ be \emph{position-balanced accuracy}.

\begin{itemize}[leftmargin=*,itemsep=0pt,topsep=0pt]
    \item \textbf{Spread}: $\max_p \acc_p - \min_p \acc_p$.
    \item \textbf{Selection share}: how often the model chooses each position; $1/k$ for an unbiased model.
    \item \textbf{Inflation}, defined three ways because ``inflates scores'' is ambiguous:
    \emph{as-published} (accuracy in the original option order minus $\bal$), which is the score people report;
    \emph{early-key} (mean of $\acc_A$ and $\acc_B$ minus $\bal$), which tests the early-position mechanism;
    and \emph{best-position} ($\max_p \acc_p - \bal$), an upper bound for a key aligned with the model's preference.
    \item \textbf{Flip rate}: share of questions answered correctly in some rotations and incorrectly in others.
    \item \textbf{Position-attributable share}: best-position inflation divided by best-position accuracy.
\end{itemize}

\para{Difficulty.}
The difficulty of a question for a model is the mean position-balanced correctness of the \emph{other six} models on that question.
It is therefore not computed from the responses being analyzed.
As a secondary measure we use each model's own confidence, taken from the separate numeric-label run.
Within difficulty strata we report a \emph{directed gap}: accuracy at the model's overall best position minus accuracy at its overall worst position.
We added this metric after planning, to avoid picking the best position separately inside each small stratum.

\subsection{Statistical analysis}

We report 95\% percentile bootstrap intervals over questions (2{,}000 resamples for per-model numbers and 1{,}000 for means across models), resampling all rotations of a question together.
Cochran's Q tests whether correctness depends on key position, and a chi-square test compares selection shares with uniform; we apply a Holm correction across all 70 model-by-condition cells.
A permutation null, which shuffles position labels within each question (300 draws), gives the spread expected from noise alone.

\para{Analyses added after seeing results.}
The permutation null, the decomposition of the as-published gain (\secref{sec:inflation}), the split of \mmlu into order-dependent and self-contained questions, the split-half replication, and the low-label-mass filter (\secref{sec:robust}) were not planned in advance.
We label them as post hoc where they appear.
